\section{Ethics: The Animal as Partner}
\label{sec:ethics}

A technology that claims to read an animal's mind carries obligations an
ordinary classifier does not. The most immediate is welfare: an owner told
``she says she loves you'' may overlook the flattened ears, the whale eye,
the lip lick that signal stress, and an owner told ``my tummy hurts'' may
over-react or dismiss genuine pain when the app stays silent. Human
observers already project their categories onto dogs, reading a
stress-driven appeasement face as guilt~\cite{sci2022explainable,horowitz2009disambiguate},
and a resonant language model amplifies the projection. Welfare-first design therefore surfaces calibrated state
estimates with uncertainty rather than sentences, never overrides
observable body signals, routes high-stakes inferences (pain, aggression,
distress) to human or veterinary review, and collects passively, without
induced stress to elicit training barks.

Honesty follows from the evaluation crisis: separate the calibrated
classifier from the narrative interface, label the latter as
interpretation, disclose that accuracy figures refer to emotion
classification, and refrain from unvalidated reverse-direction claims.
Privacy follows from the always-on home sensor: on-device inference where
possible (the DolphinGemma-on-Pixel precedent is
instructive~\cite{google2025dolphingemma}), minimal retention, explicit
household consent for cloud training, and clear deletion. The animal cannot
consent; the household must, and the design should make that consent
meaningful.

Finally, a conceptual obligation. Treating the dog as an oracle whose
``words'' we extract flattens its communicative life into a human format.
The soundboard paradigm is ethically attractive precisely because it treats
the dog as a partner in a negotiated symbol system, with the dog's behavior
as the check~\cite{bastos2024plos,bastos2024scirep}; CETI's leadership has
likewise stressed that breakthroughs in whale communication should be a
collective human decision, with responsible-interaction guidelines before
signals are sent to wild animals. This framework inherits the consensus
principles of responsible AI (transparency, justice, non-maleficence,
responsibility, privacy)~\cite{jobin2019global}, the documented gap between
principles and practice~\cite{mittelstadt2019principles,floridi2018ai4people},
the political-economy critique of hidden data and labor supply
chains~\cite{crawford2021atlas}, the risk-tiered regulatory baseline of
the EU AI Act~\cite{european2024aiact}, and the concrete engineering of
fair, safe, and human-compatible systems---disparate-impact auditing,
specification-gaming avoidance, and compatibility~\cite{barocas2016bigdata,amodei2016concrete,russell2019human}. On the animal side, the precautionary
principle for sentience warrants moral consideration under
uncertainty~\cite{birch2017precaution}, and recent work locates the edge of
sentience in exactly the graded affective capacities this survey treats as
decodable~\cite{birch2024edge}. The synthesis is a double obligation: be
honest about what the model knows, and be careful toward what the animal
feels.

For builders and reviewers, the framework condenses to six checks: report
Rung-2 metrics on public held-out data with calibration, never as
``translation accuracy''; separate and label interpretation versus
measurement; validate intent claims closed-loop and blind or call them
hypotheses; collect passively, infer on-device, minimize retention, obtain
meaningful consent; route high-stakes outputs to human or veterinary review
and never override body signals; and treat the animal as a partner, keeping
welfare as the terminal criterion. None of this slows the useful science;
all of it protects the animals and the credibility of the field.
